\chapter{Support Vector Machines and Supervised Synthesis}
\label{ch:svm}

\chapterguide%
  {Linear classifiers from \chapref{ch:classification}, model-capacity ideas from \chapref{ch:training}, and dot products and norms from \chapref{ch:linear-algebra}.}
  {find a maximum-margin line on a small example; compute hinge loss; explain what $C$ and $\gamma$ do; check the kernel trick with numbers; scale and tune an SVM inside validation; and choose among all the supervised families in this book.}

Two neighborhoods sit on either side of open ground, and the council wants to build a road between them. Many routes would keep every house on the correct side. The sensible one is the route that leaves the widest clearance, because a slightly misplaced new house is then still on the correct side. A \term{support vector machine (SVM)} chooses its boundary by the same reasoning, and only the houses at the edge of the road affect where it runs.

\section{The widest street}

\begin{utpfig}[The widest street]{The widest street. A hard-margin SVM places the boundary in the middle of the widest margin between the classes; the circled support vectors on its edges determine it.}{fig:svm-margin}
\begin{tikzpicture}[scale=0.85, font=\small]
% class + (filled)
\foreach \p in {(1,1),(2,1.2),(1,2.6),(0.8,1.8),(2.4,0.8)}
  \fill \p circle (2.4pt);
% class - (open)
\foreach \p in {(4,4),(5,3.5),(3.6,4.9),(4.6,4.7),(5.2,2.9)}
  \draw[thick, fill=UTPPaleOrange] \p circle (2.4pt);
% support vectors
\fill (1.5,2.5) circle (2.4pt);
\draw[thick, fill=UTPPaleOrange] (2.5,3.5) circle (2.4pt);
\draw[thick, fill=UTPPaleOrange] (4,2) circle (2.4pt);
\draw (1.5,2.5) circle (5.5pt);
\draw (2.5,3.5) circle (5.5pt);
\draw (4,2) circle (5.5pt);
% lines
\draw[thick] (0,5) -- (5,0);
\draw[UTPMuted,thin] (4.6,0.4) -- (5.45,-0.35) node[right=-1pt] {$\bm{w}^{\mathsf T}\bm{x}+b=0$};
\draw[dashed] (0,4) -- (4,0);
\draw[dashed] (0,6) -- (6,0);
% margin arrow
\draw[{Stealth[length=2mm]}-{Stealth[length=2mm]}, thick]
  (3.6,0.4) -- (4.6,1.4) node[pos=1, above right=0pt and 1pt] {margin};
\node[align=left] at (6.0,5.35) {circled points =\\support vectors};
\end{tikzpicture}
\end{utpfig}

\begin{workedexample}
Six points: $(1,1)$, $(2,1)$, $(1,2)$ are class $-1$, and $(3,3)$, $(4,3)$, $(3,4)$ are class $+1$ (\figref{fig:svm-toy}). The SVM's answer is $\bm{w}=(\tfrac23,\tfrac23)$ and $b=-3$. Check the score $y(\bm{w}^{\mathsf T}\bm{x}+b)$ at each point: for $(2,1)$ it is $-1\times(\tfrac23\cdot3-3)=1$, and for $(1,2)$ and $(3,3)$ it is also exactly 1. These three are the \term{support vectors}. The other three score 1.67, so they sit beyond the curb and do not affect the street at all. Delete them and the answer stays the same.

The street's width is $2/\|\bm{w}\|_2=2/(\tfrac23\sqrt2)=2.12$.
\end{workedexample}

\begin{utpfig}[The widest street on six points]{The fitted boundary and its two curbs on six points. The three circled support vectors sit exactly on the curbs. The street is 2.12 wide.}{fig:svm-toy}
\begin{tikzpicture}
\begin{axis}[logaxis, width=7.2cm, height=6.4cm, axis equal image, xmin=0, xmax=5, ymin=0, ymax=5, xlabel={$x_1$}, ylabel={$x_2$}, xtick={0,...,5}, ytick={0,...,5}]
\addplot[draw=UTPNavy,very thick,domain=0:4.5] {4.5-x};
\addplot[draw=UTPNavy,dashed,domain=0:3] {3-x};
\addplot[draw=UTPNavy,dashed,domain=1:5] {6-x};
\addplot[only marks,mark=*,mark size=2.2pt,draw=UTPBlue,fill=UTPBlue] coordinates {(1,1)(2,1)(1,2)};
\addplot[only marks,mark=square*,mark size=2.2pt,draw=UTPOrange,fill=UTPOrange] coordinates {(3,3)(4,3)(3,4)};
\addplot[only marks,mark=o,mark size=4.2pt,draw=UTPRed,thick] coordinates {(2,1)(1,2)(3,3)};
\node[font=\scriptsize\color{UTPNavy},anchor=west] at (axis cs:0.15,0.35) {class $-1$};
\node[font=\scriptsize\color{UTPNavy},anchor=east] at (axis cs:4.9,4.6) {class $+1$};
\end{axis}
\end{tikzpicture}
\end{utpfig}

\section{The hard-margin objective}

For $y_i\in\{-1,+1\}$, require correctly signed linear scores of magnitude at least one:
\[
  \min_{\bm{w},b}\ \frac{1}{2}\|\bm{w}\|_2^2
  \quad\text{subject to}\quad
  y_i(\bm{w}^{\mathsf T}\bm{x}_i+b)\geq 1 .
\]
The constraint requires the correct sign and score magnitude at least~1. The margin width is $2/\|\bm w\|_2$, so minimizing the norm maximizes that width.

\section{The soft margin: allowing exceptions}

Real classes overlap, so the SVM used in practice permits violations and charges for them:
\[
  \min_{\bm{w},b}\
  \frac{1}{2}\|\bm{w}\|_2^2
  +C\sum_{i=1}^{n}
  \max\!\left(0,\,1-y_i(\bm{w}^{\mathsf T}\bm{x}_i+b)\right).
\]
\term{Hinge loss} is zero beyond the margin and linear within it or on the wrong side. $C$ balances margin width against violations:

\begin{workedexample}
The hinge loss is $\max(0,\,1-m)$, where $m=y(\bm{w}^{\mathsf T}\bm{x}+b)$ measures how far a point is on the right side, in margin units.

\medskip
\centering\small
\begin{tabular}{@{}lcc@{}}
\toprule
\rowcolor{UTPPaleBlue!92}
\utpth{Where the point is} & \utpth{$m$} & \utpth{Hinge loss} \\
\midrule
Well beyond the curb & 2.0 & 0.0 \\
Exactly on the curb & 1.0 & 0.0 \\
Inside the street, right side & 0.4 & 0.6 \\
On the wrong side & $-0.5$ & 1.5 \\
\bottomrule
\end{tabular}
\par\raggedright
\medskip
Points beyond the curb cost nothing and so have no effect on the fit. Log loss, by contrast, keeps rewarding more confidence forever.
\end{workedexample}

\begin{itemize}
  \item \term{Large $C$}: violations cost more; weaker regularization can overfit.
  \item \term{Small $C$}: more violations are tolerated; regularization is stronger.
\end{itemize}

Both SVMs and regularized logistic regression combine loss with an $L_2$ penalty. Hinge loss becomes zero beyond the margin; cross-entropy continues rewarding confidence.

\section{Curved boundaries: mapping up}

Some datasets cannot be separated by any straight line --- but they can be after the features are transformed.

\begin{workedexample}
Let $x\in\{-2,2\}$ be class A and $x\in\{-1,0,1\}$ class B. No single threshold separates them. Map to $\phi(x)=(x,x^2)$:
\begin{utptable}
\begin{tabular}{@{}lccccc@{}}
\toprule
\rowcolor{UTPPaleBlue!92}
\utpth{$x$} & \utpth{$-2$} & \utpth{$-1$} & \utpth{$0$} & \utpth{$1$} & \utpth{$2$} \\
\midrule
$x^2$ & $4$ & $1$ & $0$ & $1$ & $4$ \\
class & A & B & B & B & A \\
\bottomrule
\end{tabular}
\end{utptable}
Class A has height 4 and class B at most 1. The line $x^2=2.5$ separates them, corresponding to two thresholds $|x|=\sqrt{2.5}$ in the original space.
\end{workedexample}

\begin{analogy}
Marbles of two colors lie on a table, with one color in a ring around the other. No straight stick can separate them. Give the table a sharp tap from below so the central marbles fly up, and a flat sheet of card can slide between the two groups. A new feature gives you the extra height.
\end{analogy}

\section{The kernel trick}

Rich feature maps can be expensive or infinite-dimensional. The SVM's equations happen to use the mapped rows only through dot products, and a \term{kernel function} computes those dot products directly:
\[
  k(\bm{x},\bm{x}')=\phi(\bm{x})^{\mathsf T}\phi(\bm{x}'),
\]
so the mapped features never need to be constructed at all. That shortcut is the \term{kernel trick}.

\begin{workedexample}
Let $\phi(\bm{v})=(v_1^2,\ \sqrt2\,v_1v_2,\ v_2^2)$, and take $\bm{x}=(1,2)$ and $\bm{z}=(3,1)$. The long way: $\phi(\bm{x})=(1,\ 2.83,\ 4)$ and $\phi(\bm{z})=(9,\ 4.24,\ 1)$, whose dot product is $9+12+4=25$. The short way: $(\bm{x}^{\mathsf T}\bm{z})^2=(3+2)^2=25$. Same answer, with no need to visit the larger space.
\end{workedexample}

Standard choices:

\begin{utptable}
\begin{tabular}{@{}lll@{}}
\toprule
\rowcolor{UTPPaleBlue!92}
\utpth{Kernel} & \utpth{Formula} & \utpth{Boundary it produces} \\
\midrule
Linear & $\bm{x}^{\mathsf T}\bm{x}'$ & Straight; fast on huge, sparse data \\
Polynomial & ${(\gamma\,\bm{x}^{\mathsf T}\bm{x}'+r)}^d$ & Curved, degree-$d$ interactions \\
RBF (Gaussian) & $\exp\!\left(-\gamma\|\bm{x}-\bm{x}'\|_2^2\right)$ & Smooth, local, very flexible \\
\bottomrule
\end{tabular}
\end{utptable}

The \term{RBF kernel} measures similarity that decays with distance. Small $\gamma$ gives broad influence and smooth boundaries; large $\gamma$ gives local, flexible boundaries. Tune $C$ and $\gamma$ together on logarithmic grids.

\begin{workedexample}
RBF similarity $\exp(-\gamma d^2)$ between two points at distance $d$:

\medskip
\centering\small
\begin{tabular}{@{}cccc@{}}
\toprule
\rowcolor{UTPPaleBlue!92}
\utpth{$\gamma$} & \utpth{$d=0.5$} & \utpth{$d=1$} & \utpth{$d=2$} \\
\midrule
0.1 & 0.975 & 0.905 & 0.670 \\
1 & 0.779 & 0.368 & 0.018 \\
10 & 0.082 & 0.000 & 0.000 \\
\bottomrule
\end{tabular}
\par\raggedright
\medskip
At $\gamma=10$ a training point has almost no influence even one unit away, so the boundary can wrap individual points.
\end{workedexample}

\begin{intuition}
RBF support vectors act as landmarks with distance-decaying influence. The kernel evaluates their similarity without explicitly constructing the feature map.
\end{intuition}

\subsection{Linear versus RBF boundaries, drawn}

A feature map can make a nonlinear separation linear in the transformed space.

\begin{utpfig}[Linear versus RBF boundaries]{Linear versus RBF boundaries. One straight line cannot put a central region in its own class; an RBF kernel draws a smooth curved boundary that isolates it.}{fig:svm-kernels}
\begin{tikzpicture}[font=\scriptsize]
\begin{scope}
  \draw[draw=UTPMuted!70,fill=white] (0,0) rectangle (4.7,3.8);
  \node[font=\small\bfseries\color{UTPNavy}] at (2.35,4.2) {Linear kernel};
  \foreach \p in {(2.0,1.7),(2.4,1.9),(2.1,2.3),(2.7,2.2),(2.55,1.55)}
    \fill[UTPBlue] \p circle (2.6pt);
  \foreach \p in {(0.65,0.65),(1.0,3.05),(3.75,0.7),(4.0,3.0),(0.55,2.0),(4.15,1.85)}
    \draw[UTPOrange,very thick] \p circle (2.9pt);
  \draw[UTPNavy,very thick] (0.35,2.9) -- (4.35,0.85);
  \node[align=center,text=UTPRed] at (2.35,-0.48)
    {one straight line cannot\\put the center in its own class};
\end{scope}

\begin{scope}[xshift=6.0cm]
  \draw[draw=UTPMuted!70,fill=white] (0,0) rectangle (4.7,3.8);
  \node[font=\small\bfseries\color{UTPNavy}] at (2.35,4.2) {RBF kernel};
  \foreach \p in {(2.0,1.7),(2.4,1.9),(2.1,2.3),(2.7,2.2),(2.55,1.55)}
    \fill[UTPBlue] \p circle (2.6pt);
  \foreach \p in {(0.65,0.65),(1.0,3.05),(3.75,0.7),(4.0,3.0),(0.55,2.0),(4.15,1.85)}
    \draw[UTPOrange,very thick] \p circle (2.9pt);
  \draw[UTPTeal,very thick] (2.35,1.95) ellipse (1.15 and 0.92);
  \node[align=center,text=UTPGreen] at (2.35,-0.48)
    {a smooth curved boundary\\can isolate the local region};
\end{scope}

\node[font=\scriptsize\color{UTPBlue}] at (4.9,4.18) {$\bullet$ class A};
\node[font=\scriptsize\color{UTPOrange}] at (6.65,4.18) {$\circ$ class B};
\end{tikzpicture}
\end{utpfig}

\begin{keyidea}
A kernel changes the representation during fitting; it does not bend an already-trained line afterward.
\end{keyidea}

\keepspace{18\baselineskip}
\section{What $C$ and $\gamma$ actually do}
\label{sec:svm-controls}

\begin{utpfig}[What C does]{What $C$ does. A small $C$ keeps a wide street and tolerates a stray point; a large $C$ narrows the street and bends to capture it, at the risk of overfitting.}{fig:svm-c}
\begin{tikzpicture}[scale=0.82, every node/.style={font=\scriptsize}]
\begin{scope}
\draw[draw=UTPMuted,fill=white] (0,0) rectangle (3.5,2.8);
\draw[UTPNavy,very thick] (0.3,0.2) -- (3.2,2.6);
\draw[UTPNavy,dashed] (0.05,0.75) -- (2.7,2.95);
\draw[UTPNavy,dashed] (0.75,-0.15) -- (3.5,2.15);
\fill[UTPBlue!12] (0.05,0.75) -- (2.7,2.95) -- (3.5,2.15) -- (0.75,-0.15) -- cycle;
\draw[UTPNavy,very thick] (0.3,0.2) -- (3.2,2.6);
\foreach \p in {(0.5,1.9),(1.0,2.4),(0.35,1.35),(1.5,2.65)} \fill[UTPBlue] \p circle (2.2pt);
\foreach \p in {(2.4,0.5),(2.9,1.0),(3.2,0.4),(1.9,0.25)} \fill[UTPOrange] \p circle (2.2pt);
\fill[UTPOrange] (1.65,1.6) circle (2.2pt);
\node[font=\small\bfseries\color{UTPNavy}] at (1.75,3.15) {small $C$: wide street};
\node[color=UTPMuted,align=center] at (1.75,-0.55) {tolerates the stray point\\simpler, may underfit};
\end{scope}
\begin{scope}[xshift=4.6cm]
\draw[draw=UTPMuted,fill=white] (0,0) rectangle (3.5,2.8);
\fill[UTPRed!10] (0.35,0.95) -- (2.15,2.9) -- (2.6,2.5) -- (0.8,0.55) -- cycle;
\draw[UTPRed,very thick,rounded corners=2pt] (0.2,0.5) -- (1.2,1.3) -- (1.45,1.85) -- (2.0,2.2) -- (2.6,2.8);
\foreach \p in {(0.5,1.9),(1.0,2.4),(0.35,1.35),(1.5,2.65)} \fill[UTPBlue] \p circle (2.2pt);
\foreach \p in {(2.4,0.5),(2.9,1.0),(3.2,0.4),(1.9,0.25)} \fill[UTPOrange] \p circle (2.2pt);
\fill[UTPOrange] (1.65,1.6) circle (2.2pt);
\draw[UTPRed,very thick] (1.4,1.45) circle (0.32);
\node[font=\small\bfseries\color{UTPNavy}] at (1.75,3.15) {large $C$: narrow street};
\node[color=UTPMuted,align=center] at (1.75,-0.55) {bends to capture it\\complex, may overfit};
\end{scope}
\end{tikzpicture}
\end{utpfig}

\begin{utptable}
\begin{tabularx}{\linewidth}{@{}>{\bfseries\leavevmode\color{UTPNavy}}m{0.10\linewidth}m{0.27\linewidth}X@{}}
\toprule
\rowcolor{UTPPaleBlue!92}
\utpth{Dial} & \utpth{Plain meaning} & \utpth{Turn it up when} \\
\midrule
$C$ & How badly you want to avoid mistakes on the training data & The model underfits; the boundary is too simple \\
$\gamma$ & How far a single training point's influence reaches & The boundary needs to be more local and wiggly \\
\bottomrule
\end{tabularx}
\end{utptable}

\begin{analogy}
$C$ is a strictness setting for a teacher marking homework. A low $C$ teacher accepts that a few students will get things wrong and teaches a general rule. A high $C$ teacher insists that every single student's answer be accommodated, and ends up teaching a rule so contorted that nobody can apply it to a new question.

$\gamma$ is the radius of each training point's voice. Low $\gamma$ means every point is heard across the whole room, producing a smooth global boundary. High $\gamma$ means each point is only heard by its immediate neighbors, producing islands around individual examples.
\end{analogy}

\begin{workedexample}
Lab~3's grid search on 15,000 personal loans, RBF kernel with balanced class weights, three-fold cross-validated ROC AUC on the training rows (\texttt{scale} means $\gamma=1/p$ on standardized data, here about 0.077 for thirteen columns):

\medskip
\centering\small
\begin{tabular}{@{}cccc@{}}
\toprule
\rowcolor{UTPPaleBlue!92}
\utpth{$C$} & \utpth{$\gamma=0.01$} & \utpth{$\gamma=$ scale} & \utpth{$\gamma=0.1$} \\
\midrule
0.1 & \textbf{0.718} & 0.714 & 0.713 \\
1 & 0.715 & 0.707 & 0.702 \\
10 & 0.714 & 0.665 & 0.647 \\
\bottomrule
\end{tabular}
\par\raggedright
\medskip
Every cell lies between 0.647 and 0.718. The worst corner pairs the narrowest kernel with the largest $C$: each loan only speaks for its own neighborhood, and the boundary bends to fit every one. The winner, $C=0.1$ with $\gamma=0.01$, is the smoothest setting in the grid, and it still keeps 11,478 of its 15,000 training loans as support vectors, 2,205 of them defaults: on noisy data most loans sit near the boundary. The best linear kernel scores 0.713 on 20,000 fresh test loans, against 0.716 for the RBF kernel, so here the curvature does almost no work.
\end{workedexample}

\section{Multiclass SVMs and support vector regression}

Multiclass reductions include \term{one-vs-rest} ($K$ classifiers) and \term{one-vs-one} (one per class pair, combined by voting).

\term{Support vector regression (SVR)} recycles the margin idea for numeric labels: build a tube of half-width $\varepsilon$ around the prediction and charge nothing for points inside it, using the $\varepsilon$-insensitive loss $\max(0,\,|y-f(\bm{x})|-\varepsilon)$. Only points on or outside the tube become support vectors.

\section{Practical notes}

\begin{itemize}
  \item \term{Scale the features.} Apply the distance-geometry rule from \secref{sec:scaling}.
  \item \term{Tune $C$ (and $\gamma$)} by cross-validation over logarithmic grids such as $10^{-2},\ldots,10^{3}$.
  \item \term{Scale limits:} kernel SVMs suit small-to-medium datasets; linear SVMs suit high-dimensional text. Large datasets may favor linear models or ensembles.
  \item \term{Probabilities are not native.} SVMs output scores and margins; if calibrated probabilities are needed, apply and validate a calibration step (\chapref{ch:evaluation}).
\end{itemize}

\recipecard{Support vector machine}{A separating line or plane, $\bm{w}^{\mathsf T}\phi(\bm{x})+b$, possibly in a kernel's richer space.}{Hinge loss plus a penalty that keeps the street wide; $C$ sets the balance.}{A convex quadratic program. Its one best answer depends only on the support vectors.}

\section{Choosing a supervised model}
\label{sec:supervised-model-choice}

\begin{utptable}
\begin{tabular}{@{}p{0.19\linewidth}p{0.36\linewidth}p{0.36\linewidth}@{}}
\toprule
\rowcolor{UTPPaleBlue!92}
\utpth{Model} & \utpth{Useful strengths} & \utpth{Important limitations} \\
\midrule
Linear regression & Interpretable, fast, the standard first baseline & Linear relationships only; sensitive to outliers \\
Ridge / lasso & Controls overfitting; lasso selects features & Still linear; needs feature scaling \\
Logistic regression & Strong probability baseline, convex training, interpretable odds & Linear boundary; calibration still needs checking \\
LDA / QDA & Fast, stable on small data & Gaussian class assumptions \\
$k$-NN & No training, flexible boundaries & Slow at prediction; scale- and dimension-sensitive \\
Naive Bayes & Very fast, strong text baseline & Independence and calibration caveats (\secref{sec:naive-bayes}) \\
Decision tree & Readable rules, no scaling, mixed data & Unstable; overfits without pruning \\
Random forest & Robust, accurate, minimal tuning & Larger, slower, less transparent \\
Gradient boosting & Often a very strong tabular benchmark & Needs careful tuning and early stopping \\
SVM & Margins and kernels; strong in high dimensions & Scaling required; slow on very large data \\
\bottomrule
\end{tabular}
\end{utptable}

Start with regularized linear models, and add complexity only for validated gains. In Lab~3 a tuned RBF SVM reaches a test ROC AUC of 0.716 on fresh loans, where $k$-NN scored 0.723, naive Bayes 0.716 and plain logistic regression 0.731 on validation. A linear SVM reaches 0.713. Here the extra complexity is not earned: the SVM is slower, its saved file holds 11,478 support vectors, and it ranks loans no better than a straight line. On data with a curved boundary the RBF kernel can win clearly; on noisy data like this, the simpler model is the one to keep.

\begin{modelcard}{Support Vector Machines - Quick Guide}
\begin{tabularx}{\linewidth}{@{}>{\bfseries\leavevmode\color{UTPNavy}}p{0.24\linewidth}X@{}}
Best when & The dataset is small to medium sized, features are informative, and a strong margin-based boundary is useful. \cr
Preprocessing & Standardize numerical features; encode categories; remove clearly irrelevant features when justified by the validated workflow in \secref{sec:feature-selection}. \cr
Main controls & $C$ for margin violations; kernel choice; $\gamma$ for RBF-kernel reach; $\varepsilon$ for SVR. \cr
Watch for & Slow training on very large datasets, sensitive scale, difficult probability interpretation, and expensive kernel tuning. \cr
\end{tabularx}
\end{modelcard}

\begin{chapterrecap}
\begin{itemize}
  \item An SVM chooses the widest street between the classes. Only the support vectors, the points on or inside its curbs, affect it.
  \item Soft margins charge hinge loss for violations, and $C$ sets how expensive those are.
  \item Mapping to richer features can make curved boundaries straight, and a kernel computes the needed dot products without visiting that space.
  \item $\gamma$ sets how far one point's influence reaches. Tune $C$ and $\gamma$ together, with scaling inside the folds.
  \item Across the supervised families, start simple and keep complexity only when validation pays for it.
\end{itemize}
\end{chapterrecap}

\begin{quickcheck}
\begin{enumerate}
  \item What is the geometric meaning of a larger separating margin in an SVM?
  \item In the six-point example, what happens to the fitted line if you delete the point $(4,3)$? Why?
  \item A point has $y(\bm{w}^{\mathsf T}\bm{x}+b)=0.7$. What is its hinge loss, and where is it?
  \item How do $C$ and, for an RBF kernel, $\gamma$ change the flexibility of the fitted decision boundary?
  \item Why should scaling and SVM tuning both happen inside the validation procedure?
\end{enumerate}
\end{quickcheck}

\begin{checkanswers}
\begin{enumerate}
  \item A wider street between the classes, which leaves more room for new points to land on the correct side. Maximizing it is the SVM's way of controlling capacity.
  \item Nothing. $(4,3)$ is not a support vector: it lies beyond the curb with a score of 1.67, so it exerts no force on the solution.
  \item $1-0.7=0.3$. The point is on the correct side but inside the street.
  \item A larger $C$ makes violations costlier, which narrows the street and allows a more complex fit. A larger $\gamma$ shortens each point's reach, which makes the boundary more local and wiggly.
  \item Margins are in feature units, so a scaler fitted on validation rows leaks, and an unscaled search tunes the wrong problem. Both belong in a pipeline refitted in every fold.
\end{enumerate}
\end{checkanswers}

\labsection{3}{Supervised learning II: model families}{sec:lab03-supervised-models}

\begin{labproject}{Neighbours, naive Bayes, trees, ensembles, and support vector machines on one workflow}
\textbf{Coverage.} Chapters~\ref{ch:knn-nb}--\ref{ch:svm}.\\
\textbf{Goal.} Compare model families under Lab~2's fixed validation workflow; inspect boundaries, tree depth, ensemble gains, and SVM tuning.

\medskip\noindent\textbf{Setup.}\par
\begin{lstlisting}[style=utppython]
from pathlib import Path

import joblib
import matplotlib.pyplot as plt
import numpy as np
import pandas as pd
from sklearn.compose import ColumnTransformer
from sklearn.ensemble import GradientBoostingRegressor, RandomForestRegressor
from sklearn.impute import SimpleImputer
from sklearn.inspection import DecisionBoundaryDisplay, permutation_importance
from sklearn.linear_model import LogisticRegression
from sklearn.metrics import accuracy_score, mean_absolute_error, mean_squared_error, roc_auc_score
from sklearn.model_selection import (
    GridSearchCV, KFold, StratifiedKFold, cross_val_score, train_test_split,
)
from sklearn.naive_bayes import GaussianNB
from sklearn.neighbors import KNeighborsClassifier
from sklearn.pipeline import Pipeline
from sklearn.preprocessing import OneHotEncoder, StandardScaler
from sklearn.svm import SVC
from sklearn.tree import DecisionTreeClassifier, export_text, plot_tree

PROJECT_ROOT = Path.cwd().parent
DATA = PROJECT_ROOT / "all_datasets"
MODELS = PROJECT_ROOT / "models"
MODELS.mkdir(exist_ok=True)
\end{lstlisting}

\medskip\noindent\textbf{Part A --- KNN versus Gaussian naive Bayes.}\par

KNN votes among standardized neighbors; Gaussian naive Bayes combines per-feature Gaussian likelihoods and class priors.

\begin{lstlisting}[style=utppython]
cancer = pd.read_csv(DATA / "breast_cancer_dataset.csv")

X_c = cancer.drop(columns=["id", "diagnosis", "Unnamed: 32"], errors="ignore")
y_c = (cancer["diagnosis"] == "M").astype(int)

Xc_train, Xc_valid, yc_train, yc_valid = train_test_split(
    X_c, y_c, test_size=0.25, random_state=42, stratify=y_c
)
\end{lstlisting}

$k$ is a hyperparameter, so it is chosen by cross-validation inside the training partition. Small $k$ follows every local quirk (low bias, high variance); large $k$ smooths everything (high bias, low variance).

\begin{lstlisting}[style=utppython]
cv = StratifiedKFold(n_splits=5, shuffle=True, random_state=42)
k_values = [1, 3, 5, 7, 11, 21]
cv_scores = {}

for k in k_values:
    candidate = Pipeline([
        ("scale", StandardScaler()),
        ("model", KNeighborsClassifier(n_neighbors=k, weights="distance")),
    ])
    scores = cross_val_score(candidate, Xc_train, yc_train, cv=cv, scoring="roc_auc")
    cv_scores[k] = scores.mean()
    print(f"k={k:2d}  cross-validated ROC AUC = {scores.mean():.3f}")

best_k = max(cv_scores, key=cv_scores.get)
print("selected k:", best_k)

plt.figure(figsize=(5, 3.2))
plt.plot(k_values, [cv_scores[k] for k in k_values], marker="o")
plt.axvline(best_k, linestyle="--", color="grey")
plt.xlabel("k (number of neighbours)")
plt.ylabel("Cross-validated ROC AUC")
plt.title("Choosing k on training data only")
plt.tight_layout()
plt.show()
\end{lstlisting}

\textbf{Inspect.} Validation improves beyond $k=1$ and is similar for $k=7$--21.

With $k$ locked, both models are compared on the validation split.

\begin{lstlisting}[style=utppython]
models = {
    "KNN": Pipeline([
        ("scale", StandardScaler()),
        ("model", KNeighborsClassifier(n_neighbors=best_k, weights="distance")),
    ]),
    "GaussianNB": Pipeline([
        ("scale", StandardScaler()),
        ("model", GaussianNB()),
    ]),
}

for name, model in models.items():
    model.fit(Xc_train, yc_train)
    pred = model.predict(Xc_valid)
    prob_c = model.predict_proba(Xc_valid)[:, 1]
    print(f"{name:10s} accuracy={accuracy_score(yc_valid, pred):.3f} "
          f"roc_auc={roc_auc_score(yc_valid, prob_c):.3f}")
\end{lstlisting}

Compare two-feature boundaries on tumor radius and texture, adding logistic regression as a reference.

\begin{lstlisting}[style=utppython]
two = ["radius_mean", "texture_mean"]
fig, axes = plt.subplots(1, 3, figsize=(12, 3.8))
for ax, (name, estimator) in zip(axes, [
    ("Logistic regression", LogisticRegression(max_iter=2000)),
    (f"KNN (k={best_k})", KNeighborsClassifier(n_neighbors=best_k, weights="distance")),
    ("Gaussian naive Bayes", GaussianNB()),
]):
    small = Pipeline([("scale", StandardScaler()), ("model", estimator)])
    small.fit(Xc_train[two], yc_train)
    DecisionBoundaryDisplay.from_estimator(small, Xc_valid[two], ax=ax, alpha=0.3,
                                           response_method="predict", cmap="coolwarm")
    ax.scatter(Xc_valid[two[0]], Xc_valid[two[1]], c=yc_valid, cmap="coolwarm",
               s=12, edgecolor="k", linewidth=0.3)
    ax.set_title(f"{name}: {accuracy_score(yc_valid, small.predict(Xc_valid[two])):.2f} accuracy")
    ax.set_xlabel("radius_mean")
    ax.set_ylabel("texture_mean")
plt.tight_layout()
plt.show()
\end{lstlisting}

\textbf{Inspect.} Logistic regression is linear; Gaussian naive Bayes curves; KNN follows local votes.

\medskip\noindent\textbf{Part B --- One readable decision tree, and what happens as it grows.}\par

Trees split on value order, so numerical preprocessing imputes missingness without scaling.

\begin{lstlisting}[style=utppython]
titanic = pd.read_csv(DATA / "titanic_dataset.csv")

features = ["Pclass", "Sex", "Age", "SibSp", "Parch", "Fare", "Embarked"]
t_num = ["Age", "SibSp", "Parch", "Fare"]
t_cat = ["Pclass", "Sex", "Embarked"]
X_t = titanic[features]
y_t = titanic["Survived"]

tree_preprocess = ColumnTransformer([
    ("num", SimpleImputer(strategy="median"), t_num),
    ("cat", Pipeline([
        ("impute", SimpleImputer(strategy="most_frequent")),
        ("onehot", OneHotEncoder(handle_unknown="ignore")),
    ]), t_cat),
])

Xt_train, Xt_valid, yt_train, yt_valid = train_test_split(
    X_t, y_t, test_size=0.25, random_state=42, stratify=y_t
)

tree = Pipeline([
    ("preprocess", tree_preprocess),
    ("tree", DecisionTreeClassifier(max_depth=4, min_samples_leaf=12, random_state=42)),
])
tree.fit(Xt_train, yt_train)
print("training accuracy:", round(accuracy_score(yt_train, tree.predict(Xt_train)), 3))
print("validation accuracy:", round(accuracy_score(yt_valid, tree.predict(Xt_valid)), 3))
\end{lstlisting}

Print and plot the tree; inspect split questions, sample counts, and classes.

\begin{lstlisting}[style=utppython]
feature_names = list(tree.named_steps["preprocess"].get_feature_names_out())
print(export_text(tree.named_steps["tree"], feature_names=feature_names, max_depth=3))

plt.figure(figsize=(12, 6))
plot_tree(tree.named_steps["tree"], feature_names=feature_names,
          class_names=["did not survive", "survived"], max_depth=2, filled=True, fontsize=8)
plt.title("Top of the fitted Titanic decision tree")
plt.tight_layout()
plt.show()
\end{lstlisting}

\textbf{Inspect.} The tree first splits on sex, then age for men and class for women.

Fit each depth and compare training with validation accuracy.

\begin{lstlisting}[style=utppython]
depths = list(range(1, 13))
train_acc, valid_acc = [], []
for depth in depths:
    t = Pipeline([
        ("preprocess", tree_preprocess),
        ("tree", DecisionTreeClassifier(max_depth=depth, random_state=42)),
    ]).fit(Xt_train, yt_train)
    train_acc.append(accuracy_score(yt_train, t.predict(Xt_train)))
    valid_acc.append(accuracy_score(yt_valid, t.predict(Xt_valid)))

plt.figure(figsize=(6, 3.6))
plt.plot(depths, train_acc, marker="o", label="training")
plt.plot(depths, valid_acc, marker="o", label="validation")
plt.xlabel("max_depth")
plt.ylabel("Accuracy")
plt.title("Deeper trees memorise; validation accuracy stops improving")
plt.legend()
plt.tight_layout()
plt.show()
\end{lstlisting}

\textbf{Inspect.} Training accuracy rises with depth while validation peaks earlier, exposing overfitting.

\medskip\noindent\textbf{Part C --- Random forest versus gradient boosting.}\par

Model $\log(1+\text{price})$ from Ames Housing's mixed-type features. Compare random forest and gradient boosting on identical three-fold splits.

\begin{lstlisting}[style=utppython]
ames = pd.read_csv(DATA / "ames_housing_dataset.csv")
X_a = ames.drop(columns=["SalePrice", "Order", "PID"], errors="ignore")   # identifiers out
y_a = np.log1p(ames["SalePrice"])

ames_preprocess = ColumnTransformer([
    ("num", SimpleImputer(strategy="median"), X_a.select_dtypes(include="number").columns),
    ("cat", Pipeline([
        ("impute", SimpleImputer(strategy="most_frequent")),
        ("onehot", OneHotEncoder(handle_unknown="ignore", sparse_output=False)),
    ]), X_a.select_dtypes(exclude="number").columns),
])
ensembles = {
    "random_forest": RandomForestRegressor(
        n_estimators=150, max_features=0.8, min_samples_leaf=2, random_state=42, n_jobs=-1
    ),
    "gradient_boosting": GradientBoostingRegressor(
        n_estimators=150, learning_rate=0.05, max_depth=3, loss="huber", random_state=42
    ),
}

cv3 = KFold(n_splits=3, shuffle=True, random_state=42)
cv_rmse = {}
for name, estimator in ensembles.items():
    pipe = Pipeline([("preprocess", ames_preprocess), ("model", estimator)])
    neg_mse = cross_val_score(pipe, X_a, y_a, cv=cv3, scoring="neg_mean_squared_error")
    cv_rmse[name] = np.sqrt(-neg_mse)
    print(f"{name:18s} CV log-RMSE = {cv_rmse[name].mean():.4f} +/- {cv_rmse[name].std():.4f}")

plt.figure(figsize=(5, 3.4))
for i, (name, values) in enumerate(cv_rmse.items()):
    plt.scatter([i] * len(values), values, s=60, label="fold scores" if i == 0 else None)
    plt.hlines(values.mean(), i - 0.2, i + 0.2, color="black", label="mean" if i == 0 else None)
plt.xticks(range(len(cv_rmse)), list(cv_rmse))
plt.ylabel("log-RMSE per fold (lower is better)")
plt.title("Fold-by-fold: is the gap bigger than the spread?")
plt.legend()
plt.tight_layout()
plt.show()
\end{lstlisting}

\textbf{Inspect.} Boosting's observed gain is modest relative to fold variation. Weigh it against tuning and runtime.

Repeat on eleven wine measurements. Permute validation columns to measure the fitted forest’s reliance on each.

\begin{lstlisting}[style=utppython]
wine = pd.read_csv(DATA / "red_wine_quality_dataset.csv")
X_w = wine.drop(columns="quality")
y_w = wine["quality"]
Xw_train, Xw_valid, yw_train, yw_valid = train_test_split(X_w, y_w, test_size=0.25, random_state=42)

wine_models = {
    "random_forest": RandomForestRegressor(n_estimators=400, min_samples_leaf=2,
                                           random_state=42, n_jobs=-1),
    "gradient_boosting": GradientBoostingRegressor(random_state=42),
}
wine_pred = {}
for name, model in wine_models.items():
    model.fit(Xw_train, yw_train)
    wine_pred[name] = model.predict(Xw_valid)
    print(f"{name:18s} MAE = {mean_absolute_error(yw_valid, wine_pred[name]):.3f}  "
          f"RMSE = {np.sqrt(mean_squared_error(yw_valid, wine_pred[name])):.3f}")

importance = permutation_importance(
    wine_models["random_forest"], Xw_valid, yw_valid, n_repeats=10, random_state=42, n_jobs=-1
)
ranked = pd.Series(importance.importances_mean, index=X_w.columns).sort_values()
spread = importance.importances_std[np.argsort(importance.importances_mean)]

fig, axes = plt.subplots(1, 2, figsize=(11, 3.8))
axes[0].barh(ranked.index, ranked.values, xerr=spread)
axes[0].set_xlabel("Increase in error when the column is shuffled")
axes[0].set_title("Permutation importance (random forest)")
axes[1].scatter(yw_valid, wine_pred["gradient_boosting"], alpha=0.5)
axes[1].plot([3, 8], [3, 8], linestyle="--", color="grey")
axes[1].set_xlabel("Observed quality")
axes[1].set_ylabel("Predicted quality")
axes[1].set_title("Gradient boosting pulls extremes to the middle")
plt.tight_layout()
plt.show()
\end{lstlisting}

\textbf{Inspect.} Alcohol and sulphates rank highly; predictions compress rare extreme ratings. Forests beat boosting here, unlike Ames.

\medskip\noindent\textbf{Part D --- Support vector machines: two dials and a kernel.}\par

Fit SVM scaling within each training fold. $C$ controls violation penalties; the kernel controls boundary shape; RBF $\gamma$ controls locality.

\begin{lstlisting}[style=utppython]
X_train, X_test, y_train, y_test = train_test_split(
    X_c, y_c, test_size=0.25, random_state=42, stratify=y_c
)

svm_pipe = Pipeline([("scale", StandardScaler()), ("model", SVC())])
search = GridSearchCV(
    svm_pipe,
    param_grid={
        "model__kernel": ["linear", "rbf"],
        "model__C": [0.1, 1, 10, 100],
        "model__gamma": ["scale", 0.01, 0.1],
    },
    cv=StratifiedKFold(5, shuffle=True, random_state=42),
    scoring="roc_auc",
    n_jobs=-1,
)
search.fit(X_train, y_train)

score = search.decision_function(X_test)
pred = search.predict(X_test)
print("best parameters:", search.best_params_)
print("cross-validated ROC AUC:", round(search.best_score_, 3))
print("test accuracy:", round(accuracy_score(y_test, pred), 3))
print("test ROC AUC:", round(roc_auc_score(y_test, score), 3))
\end{lstlisting}

Plot cross-validated AUC for each $(C,\gamma)$ pair and compare two-feature kernel boundaries.

\begin{lstlisting}[style=utppython]
results = pd.DataFrame(search.cv_results_)
rbf = results[results["param_model__kernel"] == "rbf"].copy()
rbf["gamma"] = rbf["param_model__gamma"].astype(str)
grid = rbf.pivot(index="param_model__C", columns="gamma", values="mean_test_score")

fig, axes = plt.subplots(1, 3, figsize=(13, 3.8))
im = axes[0].imshow(grid.values, cmap="viridis", aspect="auto")
axes[0].set_xticks(range(len(grid.columns)), grid.columns)
axes[0].set_yticks(range(len(grid.index)), grid.index)
axes[0].set_xlabel("gamma")
axes[0].set_ylabel("C")
for i in range(grid.shape[0]):
    for j in range(grid.shape[1]):
        axes[0].text(j, i, f"{grid.values[i, j]:.3f}", ha="center", va="center",
                     color="white", fontsize=8)
axes[0].set_title("RBF kernel: CV ROC AUC by C and gamma")
for ax, (title, estimator) in zip(axes[1:], [
    ("Linear kernel", SVC(kernel="linear", C=1)),
    ("RBF kernel (gamma = 1)", SVC(kernel="rbf", C=1, gamma=1.0)),
]):
    small = Pipeline([("scale", StandardScaler()), ("model", estimator)]).fit(X_train[two], y_train)
    DecisionBoundaryDisplay.from_estimator(small, X_test[two], ax=ax, alpha=0.3,
                                           response_method="predict", cmap="coolwarm")
    ax.scatter(X_test[two[0]], X_test[two[1]], c=y_test, cmap="coolwarm",
               s=12, edgecolor="k", linewidth=0.3)
    ax.set_title(f"{title}: {accuracy_score(y_test, small.predict(X_test[two])):.2f} accuracy")
    ax.set_xlabel("radius_mean")
    ax.set_ylabel("texture_mean")
plt.tight_layout()
plt.show()
\end{lstlisting}

\textbf{Inspect.} Most scores are similar. RBF allows curved boundaries, but these scaled classes gain little over a linear model.

Save, reload, and compare predictions with the in-memory model before sharing the artifact.

\begin{lstlisting}[style=utppython]
model_path = MODELS / "breast_cancer_svm.joblib"
joblib.dump(search.best_estimator_, model_path)
loaded = joblib.load(model_path)
same = np.array_equal(loaded.predict(X_test), pred)
print("file size (KB):", round(model_path.stat().st_size / 1024, 1))
print("reloaded model reproduces the test predictions:", same)
\end{lstlisting}

\begin{tryit}
Grow the Titanic tree with \texttt{max\_depth=None} and \texttt{min\_samples\_leaf=1}, and compare training and validation accuracy with the depth-4 tree. Then refit the SVM search with only the linear kernel and compare its test ROC AUC with the RBF result.
\end{tryit}
\end{labproject}

\begin{verification}
\begin{itemize}
  \item All model families are compared under Lab~2's fixed validation workflow.
  \item Validation accuracy, not training accuracy, chooses $k$ and the tree depth, and the overfitting point is identified.
  \item The forest-versus-boosting gain is judged against the fold-to-fold variation.
  \item The SVM's $C$, $\gamma$ and kernel are tuned on scaled inputs by cross-validation.
  \item \texttt{models/breast\_cancer\_svm.joblib} reloads and reproduces the test predictions.
\end{itemize}
\end{verification}

